\section{Analysis and Discussion}
\label{sec:analysis}

\subsection{What the compression actually buys}
The cleanest statement of \diffkv{}'s benefit is at the level of KV state, not allocator peak.
The compressed store is a \emph{pre-allocated, fixed} pool: its footprint is a constant of the
configuration ($\approx\!0.2$\,GB across $28$ layers) regardless of context, whereas the dense
full-KV cache grows linearly and crosses the store's size early and keeps going
(Fig.~\ref{fig:g1}, Table~\ref{tab:main}). This is the $\rank$-vs-$\Hkv\dhead$ gap of
\S\ref{sec:compression} made concrete: each compressed token carries $O(\rank)$ scalars
instead of $O(\Hkv\dhead)$. Because the store never grows, the runtime addresses contexts the
dense cache cannot fit, which is the qualitative capability the dense baseline lacks at the top
of our range.

\subsection{Why the allocator peak tells a different story}
The MLX allocator peak is nearly identical for the compressed runtime, the exact-decode
ablation, and the dense baseline at matched context, because the peak is set during
\emph{prefill}---a full attention over the growing native cache that all three perform---and is
dominated by weights and transient activations rather than by the retained KV state
(Fig.~\ref{fig:g2}). Compression reduces the \emph{steady-state} decode footprint, not the
prefill peak. We make this explicit rather than reporting only a process-level memory number
that would credit compression with a peak reduction it does not deliver. The practical
implication is that on a memory-constrained device the binding event is whichever is larger of
(a) the prefill peak and (b) the decode-time resident set; \diffkv{} lowers (b) decisively but
leaves (a) to be addressed by prefill-side techniques (chunking, activation release), which the
runtime already applies.

\subsection{The throughput cost}
The compressed decode is several times slower than dense decode and roughly flat across context
(Fig.~\ref{fig:g3}). The flatness is consistent with the compute profile of
\S\ref{sec:decode}: work is linear in the number of \emph{blocks} with a small $\rank$-weighted
constant, so per-token cost rises slowly. The absolute gap is an implementation property, not an
algorithmic necessity: the kernel expresses many small tensor contractions inside one compiled
graph, and on this hardware that does not reach the throughput of MLX's native fused attention
over a contiguous cache. The exact-decode ablation---identical store, but attention computed
over the retained native cache---recovers dense-like throughput, which localizes the gap to the
compressed-kernel dispatch rather than to anything the compression precludes.

\subsection{The fidelity cost}
\label{sec:fidelity}
Under the needle-in-a-haystack probe, the compressed decode does \emph{not} reliably reproduce
the buried passcode once it has left the recency window, whereas the exact-decode ablation on the
same store recovers it at every context (Table~\ref{tab:main}). The exact ablation isolates the
cause: the information needed is present in the recency window when the needle is recent, but
once a needle-bearing block is compressed at rank $16$, the truncated SVD does not preserve the
single-token detail well enough for the query to recover the exact string. This is a property of
rank-$16$ truncation, not of the runtime plumbing---the ablation proves the rest of the path is
correct. We also observe that the compressed output is \emph{non-deterministic} across runs: the
range finder uses an unseeded Gaussian, so different runs produce different (and differently
wrong) reconstructions of a compressed needle. Both are addressable (higher or
adaptive rank, an exact-token residual path, a fixed seed) and are stated as open problems rather
than worked around.

\begin{observation}
\diffkv{} trades \emph{memory and context reach} for \emph{decode throughput and single-token
recall fidelity}. The trade is favorable exactly when the cache is the binding constraint and the
task tolerates approximate long-range recall; it is unfavorable when the device has memory to
spare or the task demands verbatim recall of compressed detail.
\end{observation}
